\documentclass[10pt]{beamer}
\usetheme{metropolis}
\usepackage{graphicx}
\usepackage{listings}
\usepackage{booktabs}
\usepackage{xcolor}
\definecolor{codebg}{HTML}{F5F5F5}
\definecolor{codegreen}{HTML}{2E7D32}
\definecolor{codegray}{HTML}{757575}
\definecolor{codeblue}{HTML}{1565C0}
\lstdefinestyle{rstyle}{language=R,backgroundcolor=\color{codebg},basicstyle=\ttfamily\scriptsize,
  keywordstyle=\color{codeblue}\bfseries,stringstyle=\color{codegreen},
  commentstyle=\color{codegray}\itshape,breaklines=true,frame=single,
  rulecolor=\color{codegray!30},showstringspaces=false,tabsize=2,
  morekeywords={library,gt,tab_header,tab_spanner,fmt_number,fmt_percent,
    cols_label,tab_style,cell_fill,cell_text,cell_borders,data_color,
    tab_source_note,kable,kable_styling,row_spec,column_spec,cell_spec,
    scroll_box,pack_rows}}
\lstdefinestyle{pythonstyle}{language=Python,backgroundcolor=\color{codebg},basicstyle=\ttfamily\scriptsize,
  keywordstyle=\color{codeblue}\bfseries,stringstyle=\color{codegreen},
  commentstyle=\color{codegray}\itshape,breaklines=true,frame=single,
  rulecolor=\color{codegray!30},showstringspaces=false,tabsize=4,
  morekeywords={import,from,as,pd,GT,tab_header,fmt_number,data_color,
    tab_source_note,Styler,background_gradient,bar,format,set_caption,
    highlight_max,applymap}}
\title{Tables as Visualization}
\subtitle{Workshop 3 --- Module 9}
\author{Prof.Asoc.\ Endri Raco}
\institute{Data Visualization for Data Scientists \\ UNYT Tirana}
\date{2025/26}
\begin{document}

\begin{frame}
  \titlepage
\end{frame}

\begin{frame}{Module Roadmap}
  \begin{enumerate}
    \item Table vs chart: when to use each
    \item Five table design principles
    \item Conditional formatting: colour-encoding in cells
    \item Sparklines: Tufte's inline data graphics
    \item R: \texttt{gt} and \texttt{kableExtra}
    \item Python: \texttt{great\_tables} and \texttt{pandas Styler}
  \end{enumerate}
  \begin{exampleblock}{Learning Outcome}
    You will produce publication-quality styled tables with conditional
    colour, inline sparklines, and proper typographic design in both
    R and Python.
  \end{exampleblock}
\end{frame}

\section{Table vs Chart}

\begin{frame}{When to Use a Table}
  \begin{center}
    \includegraphics[width=0.88\textwidth]{figures_w03m09/table_vs_chart}
  \end{center}
  \begin{alertblock}{Pro-Tip}
    Tables are best when the reader needs \textbf{exact values} or
    \textbf{lookup by row/column}. Charts are best when the reader
    needs to see \textbf{patterns, trends, or comparisons}. When in
    doubt, ask: ``Does the reader need the number itself, or the
    pattern the numbers form?''
  \end{alertblock}
\end{frame}

\section{Table Design Principles}

\begin{frame}{Five Design Principles}
  \begin{center}
    \includegraphics[width=0.92\textwidth]{figures_w03m09/design_principles}
  \end{center}
\end{frame}

\begin{frame}{Plain vs Styled}
  \begin{center}
    \includegraphics[width=0.92\textwidth]{figures_w03m09/plain_vs_styled}
  \end{center}
\end{frame}

\section{Conditional Formatting}

\begin{frame}{Colour Encoding Within Cells}
  \begin{center}
    \includegraphics[width=0.78\textwidth]{figures_w03m09/conditional}
  \end{center}
  \begin{exampleblock}{Data Insight}
    Conditional formatting bridges the table--chart gap: cells are
    coloured by value, enabling pattern detection while preserving
    exact numbers. Use sequential colour (light$\rightarrow$dark) for
    magnitude; diverging colour (red--white--green) for
    deviation from a reference.
  \end{exampleblock}
\end{frame}

\begin{frame}[fragile]{Conditional Formatting in R vs Python}
  \begin{columns}[T]
    \begin{column}{0.48\textwidth}
      \textbf{R (gt)}
\begin{lstlisting}[style=rstyle]
library(gt)

df |>
  gt() |>
  tab_header(
    title = "Category Summary",
    subtitle = "Google Play Store") |>

  # Colour-encode Rating column
  data_color(
    columns = mean_rating,
    palette = "RdYlGn") |>

  # Format numbers
  fmt_number(columns = count,
    decimals = 0,
    use_seps = TRUE) |>
  fmt_percent(columns = paid_pct,
    decimals = 0) |>

  # Source note
  tab_source_note(
    "Source: Kaggle") |>

  # Striped rows
  opt_row_striping()
\end{lstlisting}
    \end{column}
    \begin{column}{0.48\textwidth}
      \textbf{Python (pandas Styler)}
\begin{lstlisting}[style=pythonstyle]
(df.style
    # Colour-encode Rating
    .background_gradient(
        subset=["mean_rating"],
        cmap="RdYlGn",
        vmin=3.8, vmax=4.4)

    # Bar chart in Count column
    .bar(subset=["count"],
        color="#BBDEFB")

    # Format
    .format({
        "count": "{:,.0f}",
        "mean_rating": "{:.2f}",
        "paid_pct": "{:.0%}"})

    # Highlight max
    .highlight_max(
        subset=["mean_rating"],
        color="#C8E6C9")

    # Caption
    .set_caption(
        "Google Play Store Summary")
)
\end{lstlisting}
    \end{column}
  \end{columns}
\end{frame}

\section{Sparklines}

\begin{frame}{Sparklines: Inline Data Graphics}
  \begin{center}
    \includegraphics[width=0.82\textwidth]{figures_w03m09/sparklines}
  \end{center}
  \begin{alertblock}{Pro-Tip}
    Tufte (2006) defined sparklines as ``data-intense, design-simple,
    word-sized graphics.'' They embed a trend directly in the table
    row, eliminating the need for a separate chart. In R:
    \texttt{gt::gt\_plt\_sparkline()} or \texttt{kableExtra} with
    inline SVG. In Python: \texttt{great\_tables} nanoplot or
    matplotlib inline rendering.
  \end{alertblock}
\end{frame}

\begin{frame}[fragile]{Sparklines in R (gt)}
\begin{lstlisting}[style=rstyle]
library(gt)

df |>
  # Nest trend data as list column
  group_by(category) |>
  summarise(
    current = last(rating),
    trend = list(rating),          # list column of vectors
    change = last(rating) - first(rating),
    .groups = "drop") |>
  gt() |>
  # Render sparklines from list column
  gt_plt_sparkline(
    column = trend,
    type = "ref_line",             # add reference line at mean
    same_limit = TRUE,             # shared y range across rows
    palette = c("#1565C0",         # line colour
                "#1565C0",         # area fill
                "#E53935",         # reference line
                "#E53935",         # point high
                "#2E7D32")) |>     # point low
  fmt_number(columns = current, decimals = 2) |>
  tab_header(title = "Rating Trends by Category")
\end{lstlisting}
\end{frame}

\begin{frame}[fragile]{Styled Tables in Python (great\_tables)}
\begin{lstlisting}[style=pythonstyle]
from great_tables import GT, loc, style

(GT(df)
    .tab_header(
        title="Category Summary",
        subtitle="Google Play Store")

    # Colour-encode
    .data_color(
        columns="mean_rating",
        palette="RdYlGn")

    # Nanoplot sparklines (GT 0.6+)
    .fmt_nanoplot(
        columns="trend",
        plot_type="line")

    # Format
    .fmt_number(columns="count",
        decimals=0, use_seps=True)
    .fmt_percent(columns="paid_pct",
        decimals=0)

    # Source note
    .tab_source_note("Source: Kaggle")
)
\end{lstlisting}
\end{frame}

\section{Ecosystem}

\begin{frame}{Table Packages: R \& Python}
  \begin{center}
    \includegraphics[width=0.92\textwidth]{figures_w03m09/ecosystem}
  \end{center}
\end{frame}

\begin{frame}[fragile]{kableExtra: Quick Tables for RMarkdown}
\begin{lstlisting}[style=rstyle]
library(kableExtra)

df |>
  kable(format = "html",
    col.names = c("Category","Count","Rating","Paid %"),
    align = c("l","r","r","r"),        # right-align numbers
    digits = c(0, 0, 2, 0)) |>
  kable_styling(
    bootstrap_options = c("striped","hover","condensed"),
    full_width = FALSE,
    font_size = 12) |>
  row_spec(0, bold = TRUE,             # bold header
    color = "white", background = "#1565C0") |>
  column_spec(2, bold = TRUE) |>       # bold Count column
  row_spec(1, bold = TRUE,             # highlight top row
    background = "#E3F2FD") |>
  scroll_box(height = "400px")         # scrollable
\end{lstlisting}
\end{frame}

\section{Complete Examples}

\begin{frame}[fragile]{R: gt Styled Summary}
  \begin{columns}[T]
    \begin{column}{0.52\textwidth}
\begin{lstlisting}[style=rstyle]
apps |>
  group_by(Category) |>
  summarise(
    Count = n(),
    `Mean Rating` = mean(Rating),
    `Paid %` = mean(Type == "Paid"),
    .groups = "drop") |>
  slice_max(Count, n = 5) |>
  gt() |>
  tab_header(
    title = "Top 5 Categories",
    subtitle = "Google Play Store") |>
  data_color(
    columns = `Mean Rating`,
    palette = "RdYlGn") |>
  fmt_number(columns = Count,
    use_seps = TRUE) |>
  fmt_percent(columns = `Paid %`,
    decimals = 0) |>
  tab_source_note("Source: Kaggle") |>
  opt_row_striping()
\end{lstlisting}
    \end{column}
    \begin{column}{0.45\textwidth}
      \includegraphics[width=\textwidth]{figures_w03m09/r_output}
    \end{column}
  \end{columns}
\end{frame}

\begin{frame}[fragile]{Python: pandas Styler}
  \begin{columns}[T]
    \begin{column}{0.52\textwidth}
\begin{lstlisting}[style=pythonstyle]
summary = (apps
    .groupby("Category")
    .agg(Count=("Rating","count"),
         Rating=("Rating","mean"),
         Share=("Rating",
           lambda x: len(x)/len(apps)))
    .nlargest(5, "Count"))

(summary.style
    .background_gradient(
        subset=["Rating"],
        cmap="RdYlGn",
        vmin=3.8, vmax=4.4)
    .bar(subset=["Count"],
        color="#BBDEFB", vmin=0)
    .format({
        "Count": "{:,.0f}",
        "Rating": "{:.2f}",
        "Share": "{:.0%}"})
    .set_caption(
        "Top 5 Categories")
    .set_table_styles([
        {"selector": "caption",
         "props": [
           ("font-weight","bold"),
           ("font-size","14px")]}])
)
\end{lstlisting}
    \end{column}
    \begin{column}{0.45\textwidth}
      \includegraphics[width=\textwidth]{figures_w03m09/py_output}
    \end{column}
  \end{columns}
\end{frame}

\section{Synthesis}

\begin{frame}{Key Takeaways}
  \begin{enumerate}
    \item Use \textbf{tables for exact values}, \textbf{charts for patterns}.
          When the number itself is the message, a table wins.
    \item Five design principles: \textbf{minimal gridlines},
          \textbf{right-align numbers}, \textbf{bold header},
          \textbf{highlight key cells}, \textbf{sort by a column}.
    \item \textbf{Conditional formatting} bridges table and chart:
          colour-encoded cells enable pattern detection while
          preserving exact values.
    \item \textbf{Sparklines} embed word-sized trend graphics
          directly in table rows (Tufte, 2006).
    \item In R: \texttt{gt} (declarative, polished) or \texttt{kableExtra}
          (quick, RMarkdown-friendly). In Python: \texttt{great\_tables}
          (gt-inspired) or \texttt{pandas Styler} (built-in).
  \end{enumerate}
\end{frame}

\begin{frame}{Essential Readings}
  \begin{itemize}
    \item Tufte, E.~R. (2006). \emph{Beautiful Evidence}, Chapter 2
          (Sparklines).
    \item Few, S. (2012). \emph{Show Me the Numbers}, Chapters 6--7
          (Table Design).
    \item Iannone, R. et al. (2024). gt: Easily Create
          Presentation-Ready Tables. \texttt{https://gt.rstudio.com}
    \item Iannone, R. (2024). great\_tables for Python.
          \texttt{https://posit-dev.github.io/great-tables/}
  \end{itemize}
  \vspace{6pt}
  \textbf{Next Module}: Lab: Chart Selection Decision Framework (W03-M10)
\end{frame}

\end{document}
